\documentclass[10pt]{article}

\usepackage[a4paper,margin=1in]{geometry}
\usepackage[T1]{fontenc}
\usepackage[utf8]{inputenc}
\usepackage{microtype}
\usepackage{booktabs}
\usepackage{array}
\usepackage{enumitem}
\usepackage{xcolor}
\usepackage{hyperref}
\usepackage{url}

\hypersetup{
  colorlinks=true,
  linkcolor=blue!55!black,
  citecolor=blue!55!black,
  urlcolor=blue!55!black
}
\urlstyle{same}
\emergencystretch=2em

\title{Agent Wiki: A Local-First Evidence-Closed Markdown Wiki for LLM Agents}
\author{Zijun Chu\\
\texttt{https://github.com/NimaChu/agent-wiki}}
\date{Technical Report, July 2026}

\begin{document}
\maketitle

\begin{abstract}
Large language model (LLM) agents increasingly perform long-running reading,
coding, and maintenance tasks, but their durable knowledge often remains trapped
in chat transcripts, proprietary note systems, vector databases, or ad hoc
project memories. This report introduces Agent Wiki, a local-first,
Markdown-native knowledge substrate designed for coding agents and human users
who want source-grounded, inspectable, and low-cost long-term knowledge
management. Agent Wiki separates immutable or lightly edited raw evidence from
agent-synthesized wiki pages, uses Obsidian-compatible links as a computable
knowledge graph, and defines a ``processed'' closure condition that is checked by
local linting tools: a raw source is processed only when its primary wiki targets
resolve, the wiki layer links back to the raw evidence, and follow-up flags are
closed. The system also treats images as evidence by mirroring remote images,
building image inventories, and allowing selected visuals to be promoted into
wiki answers. A read-only dashboard exposes global and local graph views,
evidence drill-downs, and health indicators such as broken links, orphaned wiki
pages, stale sources, and processed-gate failures. A case study on an Autodesk
FlexSim documentation corpus produced 1,097 raw source notes, 93 durable wiki
pages, and 5,261 graph edges, with zero unresolved links, zero orphaned wiki
pages, and zero processed-gate violations under the included lint checks. Agent
Wiki does not replace retrieval-augmented generation; instead, it offers a
file-native, auditable staging layer from which retrieval, dashboards, and agent
workflows can later be built.
\end{abstract}

\section{Introduction}

LLM agents are now capable of reading documentation, summarizing long sources,
editing code, and maintaining project-local context across multi-step workflows.
Yet durable agent knowledge remains awkward. Many practical systems fall into
one of three patterns: transient chat history, human-authored note vaults, or
retrieval-augmented generation (RAG) stacks backed by a database or vector
index. Each pattern is useful, but each leaves a gap for independent developers
and small teams who want a private, inspectable, source-grounded knowledge base
that agents can maintain without requiring a hosted service, a custom database,
or a heavy retrieval pipeline.

Agent Wiki addresses this gap with a deliberately simple design: files first,
Markdown first, and evidence first. The system is a repository layout, command
line interface, and optional dashboard for maintaining a local knowledge base in
which raw captures and durable wiki pages are separate layers. Raw notes preserve
bibliographic metadata, capture details, snapshots, content hashes, images, and
processing status. Wiki pages distill durable knowledge units such as concepts,
entities, APIs, workflows, comparisons, and recurring questions. Claims in wiki
pages link back to raw evidence using ordinary wiki links. Local scripts scan the
vault, resolve links, build graph statistics, surface maintenance queues, and
enforce a strict processed-source closure rule.

The guiding premise is that agent memory should be inspectable before it is
intelligent. Instead of beginning with embeddings or opaque hidden state, Agent
Wiki creates a substrate that humans and agents can both read, edit, diff,
lint, and version. This makes it possible to separate capture from synthesis,
review agent-maintained knowledge with ordinary developer tools, and preserve
the option to add retrieval or automation later.

This report makes four contributions:

\begin{enumerate}[leftmargin=1.5em]
  \item It describes a local-first architecture for agent-maintained knowledge
  bases that separates raw evidence from synthesized wiki knowledge.
  \item It defines an evidence-closure invariant for processed sources and
  implements local linting checks for unresolved links, orphaned wiki pages,
  missing source links, and processed-gate failures.
  \item It introduces an image-evidence workflow for mirroring, indexing, and
  promoting visual source material into answer-ready wiki pages.
  \item It reports an initial case study over a large technical documentation
  corpus, demonstrating that the approach scales to more than one thousand raw
  source notes while preserving graph health under local checks.
\end{enumerate}

\section{Design Goals}

Agent Wiki was designed around constraints that appear frequently in real
personal and small-team agent workflows.

\paragraph{Local ownership.}
The user should own the source captures, synthesized pages, images, snapshots,
and maintenance state as local files. A user should be able to inspect the
knowledge base with a text editor, search it with command line tools, and keep
private knowledge out of public code repositories.

\paragraph{Agent-native operation.}
The repository should be legible to coding agents. The project includes
workspace rules, templates, and root npm scripts so that agents can capture
sources, maintain the vault, repair links, and answer from the wiki without
depending on globally installed skills or a particular editor.

\paragraph{Evidence before synthesis.}
Raw source material should remain available after synthesis. The wiki layer is
not a replacement for source evidence; it is a compiled layer of durable
knowledge units whose claims point back to raw notes.

\paragraph{Low operational cost.}
The system should work without a hosted database, paid subscription, required
embedding service, or always-on dashboard. Visualization is generated and opened
only when requested.

\paragraph{Auditable maintenance.}
Agent-maintained knowledge needs health checks. Broken links, orphaned pages,
unprocessed sources, follow-up flags, and weakly connected pages should be
computable from local files.

\section{System Overview}

Agent Wiki is organized as a self-contained repository with the following major
directories:

\begin{center}
\begin{tabular}{>{\raggedright\arraybackslash}p{0.28\linewidth}p{0.62\linewidth}}
\toprule
Directory & Role \\
\midrule
\path|raw/| & source notes, snapshots, evidence metadata, image inventories \\
\path|wiki/| & durable synthesized knowledge pages \\
\path|templates/| & reusable raw and wiki page templates \\
\path|scripts/| & local command line workflows for capture, lint, search, repair \\
\path|tools/wiki-dashboard/| & optional read-only graph dashboard \\
\bottomrule
\end{tabular}
\end{center}

The main command line entry point is \path|scripts/karpathy-wiki.mjs|, exposed
through root npm scripts. The core commands are:
\path|wiki:status|, \path|wiki:lint|, \path|wiki:garden|,
\path|wiki:repair-links|, \path|wiki:search|, \path|wiki:capture|,
\path|wiki:images|, \path|dashboard|, and \path|dashboard:open|.

\subsection{Raw Evidence Layer}

The raw layer is the factual source layer. A raw note stores metadata such as
title, author, publication date, source URL, capture date, source type,
snapshot path, content hash, capture method, source quality, tags, related wiki
targets, and processing status. The body preserves captured content, explicit
images, extracted claims, candidate wiki links, and processing notes.

Raw notes are intentionally not treated as disposable summaries. They are the
evidence ledger from which multiple wiki pages can be synthesized. A single
source may update many durable concepts, and a single wiki concept may cite many
raw notes.

\subsection{Wiki Layer}

The wiki layer is the compiled knowledge layer. It favors one durable knowledge
unit per page: a concept, entity, method, API, workflow, comparison, or recurring
question. Pages use Obsidian-compatible links such as \texttt{[[Page Name]]} or
\texttt{[[Page Name|label]]}. Frontmatter may include typed relation hints from
a small supported vocabulary: \texttt{supports}, \texttt{challenges},
\texttt{related\_to}, \texttt{applies\_to}, \texttt{company\_of}, and
\texttt{product\_of}.

This design borrows from evergreen-note practices, especially the preference for
atomic, linked, evolving notes \cite{matuschakEvergreen}. The difference is that
Agent Wiki makes the maintenance procedure executable by agents and auditable by
local scripts.

\subsection{Graph and Resolver}

The scanner walks Markdown files under \texttt{raw/}, \texttt{wiki/},
\texttt{templates/}, and archives. It parses simple YAML frontmatter, extracts
body links and frontmatter links, resolves targets by id, title, basename, and
alias, and emits nodes, edges, unresolved links, typed relations, invalid
relations, incoming counts, and outgoing counts.

The resolver deliberately uses ordinary Markdown and wiki-link conventions
rather than a hidden database. The graph is therefore a derived artifact, not the
source of truth. If a generated JSON file is deleted, it can be rebuilt from the
Markdown vault.

\section{Evidence Closure}

The central maintenance invariant in Agent Wiki is evidence closure. A raw note
with status \texttt{processed} should satisfy three conditions:

\begin{enumerate}[leftmargin=1.5em]
  \item It has at least one related wiki target.
  \item Each related target resolves to an existing page.
  \item At least one resolved wiki target links back to the raw note.
\end{enumerate}

In addition, explicit follow-up flags prevent a source from being treated as
complete. The lint workflow reports processed-source issues such as
\path|missing-related|, \path|unresolved-related|,
\path|missing-wiki-backlink|, and \path|explicit-followup|. This turns a
normally subjective maintenance action---``this source has been processed''---into
a local, reviewable property.

Evidence closure is modest but important. It does not prove that every claim is
correct or that the synthesis is complete. It does, however, prevent a common
failure mode in agent-maintained knowledge bases: raw material is summarized,
marked done, and later becomes unreachable from the synthesized knowledge layer.

\section{Image Evidence}

Many useful sources are visual: documentation screenshots, UI states,
architecture diagrams, charts, workflows, and configuration examples. Agent Wiki
therefore treats images as evidence rather than decoration.

The image workflow scans raw notes and selected snapshots for Markdown image
syntax and HTML \texttt{img} tags, resolves relative URLs, deduplicates images,
mirrors remote images into \path|raw/assets/<source-note>/|, writes an
\path|image-index.json|, records local paths, content types, byte counts, and
download errors, and updates the raw note's \texttt{Images} section with a
preview table. Agents are instructed to promote only the most useful visual
evidence into wiki pages, keeping exhaustive inventories in raw notes.

This mechanism supports grounded answers that include local images when a visual
state is material to the answer. It also avoids a common degradation in web
captures: links to remote images rot or become difficult to associate with their
source context.

\section{Dashboard}

The dashboard is a read-only visualization surface. Its graph generator scans
\texttt{raw/} and \texttt{wiki/}, excludes navigation and maintenance files, and
writes \texttt{public/wiki-graph.json}. The frontend displays global and local
graph browsing, a wiki-only knowledge view by default, evidence drill-down for a
wiki page and its directly linked raw sources, search, group and status filters,
selected-node links and backlinks, tags, raw/wiki/link counts, and health
indicators for inbox, follow-up, stale, broken-link, and processed-gate queues.

The dashboard does not write to the vault. This read-only constraint keeps the
Markdown repository as the source of truth and lets visualization remain an
on-demand artifact rather than a runtime dependency.

\section{Implementation}

Agent Wiki is implemented with Node.js scripts and a Vite/React dashboard. The
repository is intentionally small. In the inspected version, the root scripts
directory contains roughly 94 KB of JavaScript source, and the dashboard source
contains a React graph surface plus CSS. No database server is required.

Table~\ref{tab:commands} summarizes the primary command line workflows.

\begin{table}[h]
\centering
\begin{tabular}{ll}
\toprule
Command & Purpose \\
\midrule
\texttt{npm run wiki:status} & report vault graph and queue statistics \\
\texttt{npm run wiki:lint} & validate unresolved links, source closure, and metadata \\
\texttt{npm run wiki:garden} & show maintenance queues and weakly connected pages \\
\texttt{npm run wiki:repair-links} & assist with link repair \\
\texttt{npm run wiki:search -- "query"} & search across titles, tags, aliases, and excerpts \\
\texttt{npm run wiki:capture} & create a raw source note with metadata and snapshots \\
\texttt{npm run wiki:images} & extract, mirror, and index image evidence \\
\texttt{npm run dashboard:open} & refresh graph, start local frontend, open browser \\
\bottomrule
\end{tabular}
\caption{Primary local workflows exposed as npm scripts.}
\label{tab:commands}
\end{table}

\section{Case Study: FlexSim Documentation Vault}

To evaluate the system as a practical maintenance substrate, Agent Wiki was used
to ingest and synthesize a large technical documentation corpus for Autodesk
FlexSim. The vault health commands were run on July 7, 2026. Table~\ref{tab:case}
shows the resulting graph statistics.

\begin{table}[h]
\centering
\begin{tabular}{lr}
\toprule
Metric & Value \\
\midrule
Total nodes & 1,201 \\
Graph edges & 5,261 \\
Raw source notes & 1,097 \\
Wiki pages & 93 \\
Inbox raw sources & 0 \\
Processed raw sources & 1,096 \\
Needs follow-up & 0 \\
Stale sources & 0 \\
Unresolved links & 0 \\
Invalid relation hints & 0 \\
Orphaned wiki pages & 0 \\
Processed-source issues & 0 \\
\bottomrule
\end{tabular}
\caption{Vault health for the FlexSim case study. Metrics are produced by the
included local status and lint commands.}
\label{tab:case}
\end{table}

The case study suggests that a file-native approach can remain mechanically
healthy over a documentation-scale corpus. The strongest result is not the node
count alone, but the absence of unresolved links, orphaned wiki pages, and
processed-gate violations after agent-assisted maintenance. These metrics show
that the raw evidence layer and wiki synthesis layer can be kept connected by
local checks.

\section{Relationship to RAG and Agent Memory}

RAG combines parametric language models with external non-parametric memory to
improve knowledge-intensive generation \cite{lewis2020rag}. Agent Wiki is
complementary to RAG rather than a replacement for it. RAG typically optimizes
retrieval at answer time; Agent Wiki optimizes source preservation, synthesis,
inspection, and maintenance before retrieval. A clean Agent Wiki vault can later
serve as a high-quality corpus for embeddings, search, or retrieval pipelines.

Agent memory systems such as MemGPT study how LLMs can manage context beyond a
fixed context window using memory tiers and control flow \cite{packer2023memgpt}.
Agent Wiki focuses on the external substrate that an agent can maintain: ordinary
files, explicit evidence links, and local health checks. It does not assume that
the LLM owns or hides the memory representation.

Personal AI systems such as Khoj provide search and chat over user documents and
support Markdown, PDFs, and note tools \cite{khojDocs}. Obsidian provides a
widely used local Markdown note environment with links and graph visualization
\cite{obsidian}. Agent Wiki is narrower and more procedural: it packages a vault
layout, agent instructions, capture templates, evidence-closure rules, image
evidence handling, and lintable maintenance workflows specifically for
agent-maintained knowledge bases.

\section{Limitations}

Agent Wiki currently emphasizes auditable structure over semantic automation.
The search command is lightweight and does not rank with embeddings. The system
does not automatically prove that wiki claims faithfully summarize raw sources;
it only makes source links and maintenance state inspectable. The YAML parser is
purposefully simple and may not support every frontmatter convention used by
larger static-site or note-taking ecosystems. The dashboard is read-only and
local. Finally, the case study is a single-developer, single-corpus deployment;
future work should evaluate multi-user workflows, cross-agent consistency, and
answer quality compared with chat-only and RAG-first baselines.

\section{Future Work}

Several extensions follow naturally from the current design. First, the vault
could expose a stable Model Context Protocol server so agents can query graph
state and source evidence through structured APIs while preserving Markdown as
the source of truth. Second, the system could add optional embedding indexes that
are built from the wiki layer, the raw layer, or both, allowing retrieval to be
compared against graph-guided search. Third, evidence closure could be extended
from page-level links to claim-level citations. Fourth, image evidence could be
augmented with OCR, diagram captions, and visual similarity search. Fifth,
dashboard interactions could create suggested maintenance actions while keeping
write operations explicit and reviewable.

\section{Conclusion}

Agent Wiki demonstrates that a useful LLM-agent memory substrate can begin as a
local Markdown repository rather than a hosted service or vector database. By
separating raw evidence from synthesized wiki knowledge, enforcing a processed
closure rule, indexing images as evidence, and deriving graph health from plain
files, the system offers a pragmatic bridge between personal knowledge
management, source-grounded agent answers, and future retrieval pipelines. The
initial FlexSim case study shows that the approach can support more than one
thousand raw sources while preserving link and evidence health under local lint
checks. The broader claim is simple: before agent knowledge becomes more
automated, it should become more inspectable.

\section*{Availability}

The Agent Wiki source code is available at
\url{https://github.com/NimaChu/agent-wiki}. The current report describes the
repository state inspected on July 7, 2026.

\begin{thebibliography}{9}

\bibitem{lewis2020rag}
Patrick Lewis, Ethan Perez, Aleksandra Piktus, Fabio Petroni, Vladimir
Karpukhin, Naman Goyal, Heinrich Kuttler, Mike Lewis, Wen-tau Yih, Tim
Rocktaschel, Sebastian Riedel, and Douwe Kiela.
\newblock Retrieval-augmented generation for knowledge-intensive NLP tasks.
\newblock In \emph{Advances in Neural Information Processing Systems}, 2020.
\newblock \url{https://arxiv.org/abs/2005.11401}.

\bibitem{packer2023memgpt}
Charles Packer, Sarah Wooders, Kevin Lin, Vivian Fang, Shishir G. Patil, Ion
Stoica, and Joseph E. Gonzalez.
\newblock MemGPT: Towards LLMs as operating systems.
\newblock arXiv:2310.08560, 2023.
\newblock \url{https://arxiv.org/abs/2310.08560}.

\bibitem{matuschakEvergreen}
Andy Matuschak.
\newblock Evergreen notes.
\newblock \url{https://notes.andymatuschak.org/}.

\bibitem{obsidian}
Obsidian.
\newblock Obsidian: Sharpen your thinking.
\newblock \url{https://obsidian.md/}.

\bibitem{khojDocs}
Khoj AI.
\newblock Khoj documentation.
\newblock \url{https://docs.khoj.dev/}.

\bibitem{notebooklm}
Google.
\newblock NotebookLM.
\newblock \url{https://notebooklm.google/}.

\end{thebibliography}

\end{document}
